\documentclass[11pt]{article}
\usepackage[margin=1in]{geometry}
\usepackage{hyperref}

\begin{document}

\begin{center}
\textbf{\Large Cover Letter for JMLR Submission}
\end{center}

\vspace{1em}

\noindent \textbf{Date:} February 26, 2026

\vspace{1em}

\noindent \textbf{To:} JMLR Editorial Board

\vspace{1em}

\noindent \textbf{Re:} Submission of ``Engagement as Entanglement: Variance Signatures of Bidirectional Context Coupling in Large Language Models''

\vspace{1em}

\noindent Dear Editors,

\vspace{0.5em}

I am pleased to submit the manuscript titled ``Engagement as Entanglement: Variance Signatures of Bidirectional Context Coupling in Large Language Models'' for consideration for publication in the Journal of Machine Learning Research.

\vspace{1em}

\noindent \textbf{1. Summary of Contribution}

This paper introduces an entanglement framework that provides mechanistic explanation for the ``Lost in Conversation'' phenomenon documented by Laban et al.\ (2025), who showed that LLMs exhibit 112\% increased unreliability across 200,000+ multi-turn conversations. We demonstrate that context sensitivity ($\Delta$RCI) tracks variance reduction in response embeddings ($r = 0.76$, $p = 2.37 \times 10^{-68}$, $N = 360$), revealing bidirectional context coupling: convergent entanglement (context stabilizes) and divergent entanglement (context destabilizes). The Entanglement Stability Index (ESI) predicts which models will fail in multi-turn settings.

\vspace{1em}

\noindent \textbf{2. Prior Publications}

Earlier versions of foundational concepts appear in the following preprints:
\begin{itemize}
    \item Laxman (2026a). ``Context Curves Behavior: Measuring AI Relational Dynamics with $\Delta$RCI.'' Preprints.org\\
    DOI: 10.20944/preprints202601.1881.v2
    \item Laxman (2026b). ``Standardized Context Sensitivity Benchmark.'' Preprints.org\\
    DOI: 10.20944/preprints202602.1114.v2
\end{itemize}

\textbf{Differences from prior work:} The current submission introduces the entanglement framework, variance analysis, bidirectional coupling theory, and the ESI metric---none of which appear in the preprints. The preprints establish the $\Delta$RCI metric; this paper explains its mechanistic basis. This represents substantial ($>$50\%) new contribution.

\vspace{1em}

\noindent \textbf{3. Co-author Consent}

I am the sole author of this manuscript.

\vspace{1em}

\noindent \textbf{4. Conflicts of Interest}

I have no conflicts of interest with any JMLR action editors or reviewers.

\vspace{1em}

\noindent \textbf{5. Suggested Action Editors}

Based on research alignment with the current JMLR Action Editor list, I suggest the following action editors:
\begin{enumerate}
    \item Kai-Wei Chang (UCLA) --- Large Language Models, Trustworthy NLP, Vision-Language Models
    \item Qiang Liu (University of Texas at Austin) --- Deep learning, large language models, generative models
    \item Quanquan Gu (UCLA) --- LLMs, deep learning theory, optimization
\end{enumerate}

\vspace{1em}

\noindent \textbf{6. Suggested Reviewers}

\begin{enumerate}
    \item Philippe Laban (Salesforce AI) --- author of ``Lost in Conversation'' paper
    \item Nelson Liu (Stanford) --- author of ``Lost in the Middle'' paper
    \item Nils Reimers (Hugging Face) --- sentence embeddings expert
    \item Tatsunori Hashimoto (Stanford) --- LLM evaluation
    \item Sara Hooker (Cohere) --- model evaluation and safety
\end{enumerate}

\vspace{1em}

\noindent \textbf{7. Keywords}

Large language models, context sensitivity, multi-turn conversation, variance analysis, AI evaluation, model stability, entanglement

\vspace{1em}

\noindent \textbf{8. Data and Code}

All data and code are publicly available at: \url{https://github.com/LaxmanNandi/MCH-Research}

\vspace{2em}

\noindent Thank you for considering this submission.

\vspace{1em}

\noindent Sincerely,

\vspace{1em}

\noindent Laxman M M\\
Primary Health Centre Manchi\\
Bantwal Taluk, Dakshina Kannada\\
Karnataka 574231, India\\
Email: barlax5377@gmail.com\\
ORCID: 0009-0009-0405-6531

\end{document}
